% ============================================================
%  Experiment Design + Paper Tables — Weighted CAFHT
%  Generated 2026-04-24  |  Target: 1 - alpha = 0.90
%
%  Required packages:
%    \usepackage{booktabs}
%    \usepackage{multirow}
%    \usepackage{siunitx}   % for S column alignment (optional)
%    \usepackage{xcolor}
%    \usepackage{colortbl}
% ============================================================


% ============================================================
%  SECTION: Experiment Design
% ============================================================

\section{Experiments}
\label{sec:experiments}

We evaluate Weighted CAFHT against two ablations on synthetic time series, S\&P~500 intraday returns, and medical data. All experiments set $\alpha = 0.10$ (target
coverage $1 - \alpha = 0.90$). Three algorithm variants are compared
throughout:
%
\begin{itemize}[nosep]
  \item \textbf{Weighted CAFHT}: full method — likelihood-ratio (LR) covariate
        reweighting combined with per-series ACI $\alpha_t$ updating.
  \item \textbf{Uniform weights}: LR weighting disabled (uniform calibration
        weights); per-series ACI update retained. Isolates the contribution of
        adaptive coverage correction.
  \item \textbf{Zero $\gamma$}: LR weighting enabled; ACI step size fixed at
        $\gamma = 0$ (no $\alpha_t$ update). Isolates the contribution of
        covariate reweighting.
\end{itemize}

\subsection{Synthetic Data}
\label{sec:experiments:synthetic}

\paragraph{Data-generating processes.}
We consider two synthetic DGPs, each generating $N_{\mathrm{train}} = 600$
training series, $N_{\mathrm{cal}} = 600$ calibration series, and
$N_{\mathrm{test}} = 300$ test series of length $T = 20$.

\textit{Static-X}: Each series is driven by a scalar covariate $X^{(i)}$
drawn once at series creation time.  Training and calibration covariates
follow $X \sim \mathrm{Pois}(\lambda)$ with $\lambda = 1$; under the
\emph{static shift} condition the test covariates are redrawn from
$\mathrm{Pois}(\tilde\lambda)$ with $\tilde\lambda = 2$.  Conditional on $X$,
the response follows an AR(1) process:
%
\[
  Y_t = \rho \, Y_{t-1} + \beta X + \delta t + \varepsilon_t, \quad
  \varepsilon_t \overset{\text{i.i.d.}}{\sim} \mathcal{N}(0, \sigma^2),
\]
%
with $\rho = 0.7$, $\beta = 1.0$, $\delta = 0$ (no deterministic drift), and
$\sigma = 0.2$.  Initial values are shared between train and test to hold
$Y_0$ fixed under shift.

\textit{Dynamic-X}: The covariate itself evolves as an AR(1):
$X_{t+1} = \rho_X X_t + \kappa_X t + \eta_t$, $\eta_t \sim \mathcal{N}(0,
\sigma_X^2)$.  Under the \emph{dynamic shift} condition the entire covariate
path is regenerated under shifted parameters $(\tilde\rho_X, \tilde\kappa_X,
\tilde\sigma_X)$, while $Y_0$ is held fixed.

\paragraph{Prediction model.}
At each time step $t$, an AR(1) model is fitted by OLS on all training
trajectories up to step $t$, producing point forecasts $\hat Y_{t+1}^{(i)}$.
The nonconformity score is the absolute residual
$s_t^{(i)} = |Y_t^{(i)} - \hat Y_t^{(i)}|$.

\paragraph{Likelihood-ratio estimation.}
For each time step the test set is split uniformly at random into two halves
$H_1, H_2$.  A logistic classifier with balanced class weights is fitted on
$\mathcal{D}_{\mathrm{train}} \cup H_2$ (labels $0/1$) and used to score
$H_1$, and vice versa.  This cross-half procedure prevents leakage of
deployment-distribution information into the calibration weights.
Estimated density ratios $\hat w(z) \propto \hat p(z) / (1 - \hat p(z))$ are
clipped at $5\times$ the mean weight and renormalized to sum to one.

\paragraph{ACI update and $\gamma$ selection.}
Each test series maintains its own coverage level $\alpha_t^{(i)}$, updated
after each observation via
%
\[
  \alpha_{t+1}^{(i)} = \alpha_t^{(i)}
    + \gamma \bigl(\alpha - \mathbf{1}\bigl[Y_{t+1}^{(i)} \notin
      \hat C_{t+1}^{(i)}\bigr]\bigr),
\]
%
clipped to $(10^{-6},\, 1 - 10^{-6})$.  The step size $\gamma$ is selected
every 10 steps by a held-out simulation on a three-way split of the training
data; the search grid is $\Gamma = \{0.001, 0.005, 0.01, 0.05, 0.1\}$.

\paragraph{Evaluation.}
Each configuration is replicated over 30 independent seeds (base seed 1000).
We report pooled empirical coverage (fraction of covered intervals across all
series and time steps), the mean per-seed absolute deviation
$\overline{|\Delta|} = \frac{1}{30}\sum_{s=1}^{30}
|\hat{\mathrm{cov}}_s - (1-\alpha)|$ from the target, and mean interval
width.

% -------------------------------------------------------

\subsection{Financial Data}
\label{sec:experiments:finance}

\paragraph{Data.}
We use daily S\&P~500 intraday return data from 2024 collected via
\texttt{yfinance}.  The response variable is
$Y_t = \mathrm{Close}_t / \mathrm{Open}_t - 1$ for each ticker on each
trading day.  Four covariates are constructed for each ticker-day pair:
%
\begin{itemize}[nosep]
  \item \textbf{OvernightGapReturn}: overnight return (not lagged).
  \item \textbf{Above52wLowReturn}: distance above the 52-week low, computed
        from pre-sample data to prevent lookahead bias (not lagged).
  \item \textbf{TurnoverRatio} (lagged one day).
  \item \textbf{DailyRangeReturn} (lagged one day).
\end{itemize}

\paragraph{Covariate shift design.}
We test two sector-based shift conditions.  In the \emph{Technology} setting,
all Technology-sector tickers form the test set and the remaining $\sim\!480$
tickers provide the training and calibration pool.  In the \emph{Utilities}
setting, Utilities-sector tickers form the test set.  Sector membership is
fixed by GICS classification; the resulting covariate distributions differ
meaningfully from the pool (maximum Kolmogorov–Smirnov statistic $\approx
0.49$ for Technology, $\approx 0.75$ for Utilities, driven primarily by
beta-to-market features).  A \emph{Mixed} null baseline randomly draws 15\%
of all tickers as the test set, inducing no systematic sector shift.

\paragraph{Rolling windows and train/cal split.}
Experiments are run over 13 rolling two-month windows covering
January--March 2024.  Within each window, the non-test ticker pool is split
50/50 into training and calibration sets (seed 42).  The Technology and
Utilities conditions each yield 13 independent coverage observations; the
Mixed condition uses a single window.

\paragraph{Prediction model.}
At each time step $t$, an OLS regression of $Y$ on the four covariates is
fitted pooled across all training tickers and all time steps up to $t$,
producing ticker-level forecasts $\hat Y_{t+1}^{(i)}$.  The nonconformity
score is $s_t^{(i)} = |Y_t^{(i)} - \hat Y_t^{(i)}|$.

\paragraph{Likelihood-ratio estimation.}
The LR featurizer represents each ticker at time $t$ by: the mean, standard
deviation, and AR(1) coefficient of $Y$ over a rolling window of 10 steps,
together with the temporal mean and standard deviation of each of the four
covariates over the full prefix (11 features total).  As in the synthetic
setting, a cross-half logistic classifier estimates density ratios with
balanced class weights and $5\times$-mean clipping.

\paragraph{ACI update and $\gamma$ selection.}
The per-series ACI update is identical to the synthetic setting.  The search
grid is $\Gamma = \{0.001, 0.005, 0.01, 0.05\}$ for Technology and
$\Gamma = \{0.001, 0.005, 0.01, 0.05, 0.1\}$ for Utilities; $\gamma$ is
reselected every 10 steps.

\paragraph{Evaluation.}
For Technology and Utilities we report the mean per-window coverage and the
mean per-window absolute deviation $\overline{|\Delta|}$ from the target
$1-\alpha = 0.90$ averaged across the 13 windows.  Mean interval width is
reported in intraday return units.  The Mixed null baseline is reported for a
single window.

% ============================================================
%  TABLE 1: Synthetic Results
% ============================================================

\begin{table}[ht]
\centering
\setlength{\tabcolsep}{8pt}
\renewcommand{\arraystretch}{1.15}
\caption{%
  Synthetic results ($T=20$, $\alpha=0.10$, 30 seeds, target coverage $= 0.90$).
  \textit{Coverage}: pooled empirical coverage across all series and time steps.
  $\overline{|\Delta|}$: mean per-seed absolute deviation from $1-\alpha$.
  \textit{Width}: mean prediction interval width.
}
\label{tab:synthetic}
\begin{tabular}{@{} l l S[table-format=1.3] S[table-format=1.3] S[table-format=1.3] @{}}
\toprule
\textbf{Data condition} & \textbf{Algorithm}
  & \textbf{Coverage} & $\boldsymbol{\overline{|\Delta|}}$ & \textbf{Width} \\
\midrule
\multirow{3}{*}{No shift}
  & Weighted CAFHT      & 0.901 & 0.001 & 1.395 \\
  & Uniform weights     & 0.901 & 0.001 & 1.394 \\
  & Zero $\gamma$       & 0.901 & 0.001 & 1.395 \\
\cmidrule(l){1-5}
\multirow{3}{*}{Static shift}
  & Weighted CAFHT      & 0.891 & 0.009 & 1.501 \\
  & Uniform weights     & 0.869 & 0.031 & 1.445 \\
  & Zero $\gamma$       & 0.890 & 0.010 & 1.498 \\
\cmidrule(l){1-5}
\multirow{3}{*}{Dynamic shift}
  & Weighted CAFHT      & 0.841 & 0.059 & 1.339 \\
  & Uniform weights     & 0.836 & 0.064 & 1.323 \\
  & Zero $\gamma$       & 0.838 & 0.062 & 1.332 \\
\bottomrule
\end{tabular}
\end{table}


% ============================================================
%  TABLE 2: Finance Results
% ============================================================

\begin{table}[ht]
\centering
\setlength{\tabcolsep}{8pt}
\renewcommand{\arraystretch}{1.15}
\caption{%
  Finance results (S\&P~500 intraday returns, $\alpha=0.10$, seed~42, target
  coverage $= 0.90$).
  Technology and Utilities: mean over 13 rolling windows.
  Mixed (null baseline): single window, 15\% random test draw.
  $\overline{|\Delta|}$: mean per-window absolute deviation from $1-\alpha$.
  \textit{Width}: mean interval width (intraday return units).
  Utilities uses $\gamma$-grid $\{0.001, 0.005, 0.01, 0.05, 0.1\}$;
  Technology uses $\{0.001, 0.005, 0.01, 0.05\}$.
}
\label{tab:finance}
\begin{tabular}{@{} l l S[table-format=1.3] S[table-format=1.3] S[table-format=1.4] @{}}
\toprule
\textbf{Sector} & \textbf{Algorithm}
  & \textbf{Coverage} & $\boldsymbol{\overline{|\Delta|}}$ & \textbf{Width} \\
\midrule
\multirow{3}{*}{Technology}
  & Weighted CAFHT      & 0.888 & 0.017 & 0.0505 \\
  & Uniform weights     & 0.880 & 0.024 & 0.0493 \\
  & Zero $\gamma$       & 0.872 & 0.029 & 0.0465 \\
\cmidrule(l){1-5}
\multirow{3}{*}{Utilities}
  & Weighted CAFHT      & 0.923 & 0.026 & 0.0451 \\
  & Uniform weights     & 0.924 & 0.027 & 0.0454 \\
  & Zero $\gamma$       & 0.932 & 0.034 & 0.0447 \\
\cmidrule(l){1-5}
\multirow{3}{*}{Mixed (null)}
  & Weighted CAFHT      & 0.911 & 0.011 & 0.0414 \\
  & Uniform weights     & 0.911 & 0.011 & 0.0413 \\
  & Zero $\gamma$       & 0.907 & 0.007 & 0.0417 \\
\bottomrule
\end{tabular}
\end{table}
